\documentclass[12pt,a4paper]{amsart}
\usepackage{amscd,amssymb,amsopn,amsmath,amsthm,graphics,amsfonts,enumerate,verbatim,calc}
%\usepackage[dvips]{graphicx}
%\usepackage[colorlinks=true,linkcolor=blue,citecolor=blue]{hyperref}
%\usepackage{showlabels}
%\input xy
%\xyoption{all}
%\pagestyle{empty}
\textwidth=16cm \textheight=21.2cm \topmargin=0.5cm
\oddsidemargin=0.8cm \evensidemargin=0.8cm \headheight=15pt
\headsep=1cm \numberwithin{equation}{section}
\hyphenation{semi-stable} \emergencystretch=11pt
\setcounter{page}{10}
\newtheorem{theorem}{Theorem}[section]
\newtheorem{proposition}[theorem]{Proposition}
\newtheorem{lemma}[theorem]{Lemma}
\newtheorem{corollary}[theorem]{Corollary}
\newtheorem{remark}[theorem]{Remark}
\newtheorem{example}[theorem]{Example}
\newtheorem{definition}[theorem]{Definition}
\numberwithin{equation}{section}
\begin{document}
%\pagenumbering{gobble}
\title{Distance without Flatness: Metric-Driven Clustering on Spheres, Tori
and Hyperbolic Spaces}
\author{Karthik Mattu, Adit Dhall and Thejas Nagesh Gowda}
\address{Data and Predictive Analytics Center, School of Mathematics and Statistics, College of Science, Rochester Institute of Technology,
Rochester, NY, USA}
\email{ km5503@g.rit.edu}
\address{Data and Predictive Analytics Center, School of Mathematics and Statistics, College of Science, Rochester Institute of Technology,
Rochester, NY, USA}
\email{ ad6449@g.rit.edu}
\address{Data and Predictive Analytics Center, School of Mathematics and Statistics, College of Science, Rochester Institute of Technology,
Rochester, NY, USA}
\email{ tl6153@g.rit.edu}
\maketitle

\noindent The K-means algorithm \cite{L} minimizes the within cluster sum of squared
Euclidean distances, and its centroid update relies on the linear structure
of $\mathbb{R}^{n}$. On curved surfaces this assumption breaks and squared
Euclidean cost no longer reflects the intrinsic separation between points. Here we revisit
the algorithm under three non Euclidean geometries and study how the choice
of metric reshapes the resulting clusters. On the sphere $S^{2}$ we use the
Haversine distance \cite{MJ}, on the flat torus $\mathbb{T}^{2}$ the periodic
Euclidean metric obtained by reducing coordinate differences modulo $2\pi$
along each angular direction \cite{dC}, and on the hyperbolic Poincar\'e disk
$\mathbb{D}$ we use the Poincar\'e metric induced by the conformal factor
$2/(1-\|x\|^{2})$ on the open unit disk \cite{NK}.

We view this work as a first step toward systematic clustering on non Euclidean
spaces. The three surfaces studied here represent positive curvature, flat
periodicity, and negative curvature in their simplest forms, echoing the role
that curvature plays in shaping geometric analysis on Riemannian configuration
spaces more broadly \cite{BW}, and the framework extends naturally to product
manifolds, Stiefel and Grassmann manifolds, spaces of symmetric positive
definite matrices, and further Riemannian settings encountered in modern data
analysis \cite{P}.

\begin{thebibliography}{99}
\bibitem {L} S. Lloyd, {\it Least squares quantization in PCM}, IEEE Trans.
Inform. Theory {\bf 28} (1982), 129--137.
\bibitem {MJ} K. V. Mardia and P. E. Jupp, {\it Directional Statistics},
Wiley, Chichester, 2000.
\bibitem {dC} M. P. do Carmo, {\it Riemannian Geometry}, Birkh\"auser,
Boston, 1992.
\bibitem {NK} M. Nickel and D. Kiela, {\it Poincar\'e embeddings for learning
hierarchical representations}, Advances in Neural Information Processing
Systems {\bf 30} (2017), 6338--6347.
\bibitem {BW} M. Barbosu and T. Wiandt, {\it On the Riemannian geometry of
the planetary three-body problem}, Romanian Astron. J. {\bf 29} (2019),
149--156.
\bibitem {P} X. Pennec, {\it Intrinsic statistics on Riemannian manifolds.
Basic tools for geometric measurements}, J. Math. Imaging Vis. {\bf 25}
(2006), 127--154.
\end{thebibliography}
\end{document}
